\subsection{Several letters and a linear program}
\label{subsec:cycle-jets}
\hypertarget{source-cyclic-jets}{}

For a full Frobenius orbit the substitutions at successive positions are linked. The exponents contributing to one derivative must therefore satisfy constraints around the entire cycle.

Index the \(d\) matrix variables \(Z_s\) by \(s\in\mathbb Z/d\mathbb Z\). At position \(s\) let the differential operator be a product of minors of the row-framed
matrix \(Y_s\partial_{Z_s}\). A minor of height \(h\) uses the
first \(h\) rows and a specified list of columns. Write \(A_s\)
for the number of minors, \(A'_s\) for the number of height at
least two, and \(H_s\) for the sum of their heights. Choose a set
\(T\subseteq\mathbb Z/d\) of letters at which to treat the first
row separately. Set
\[
 D_s=H_s-{\bf1}_{s\in T}A_s,\qquad
 C_s=\begin{cases} A'_s&s\in T,\\ A_s&s\notin T.\end{cases}
\]
For nonnegative real capacities \(D_s\) allow a variable
\(z_{sj}\ge0\) for each odd \(1\le j<d\).

The variable \(z_{sj}\) uses \(p^jz_{sj}\) units at \(s\) and \(z_{sj}\) at \(s+j\); the sum of all uses
at \(s\) is bounded by \(D_s\). Define
\[
 {\cal F}_p(D)=\gamma\sum_sD_s-\max\sum_{s,j}z_{sj}.
\]
By finite linear programming duality there are numbers
\(0\le b_s\le1\) with
\begin{equation}\label{eq:flow-dual}
 p^jb_s+b_{s+j}\ge1\quad(j>0\text{ odd}),\qquad
 {\cal F}_p(D)=\sum_sD_s(\gamma-b_s).
\end{equation}
Indeed the primal feasible set is nonempty and compact. Its
dual minimizes \(\sum D_s b_s\) over nonnegative \(b_s\)
with the indicated inequalities for \(j<d\). Capping each
\(b_s\) by one preserves feasibility. For any other positive
odd \(j\), reduce \(j\) modulo the even integer \(d\) and use
nonnegativity.

For \(s\in T\) put \(h_s=\min\{j>0:j\ {\rm odd},\ s+j\in T\}\),
with minimum \(\infty\) allowed. Set
\[
 \begin{split}
 T_s={}&\sum_{\substack{0<j<h_s\\ j\ {\rm odd}}}
                 (A_sp^{-j}-C_{s+j})_+
       +\sum_{\substack{j>h_s\\ j\ {\rm odd}}}A_sp^{-j},\\
 W_s={}&\max\{C_{s-j}: j>0\ {\rm odd},\ p^j\le D_{s-j}\},
 \end{split}
\]
where \(t_+=\max(t,0)\) and an empty maximum is zero. For each
destination \(t\) form the list of pairs \((s,j)\), where
\(s\in T\), \(j\) is a signed odd integer, \(s+j=t\), and
\[
       j>0\ \Longrightarrow\ p^j\le A_s,\qquad
       j<0\ \Longrightarrow\ p^{-j}\le D_t .
\]
Denote its length by \(J_t\), and set
\(a_*=\max(1,\max_{s\in T}A_s,\max_t C_t)\).

\begin{proposition}\label{prop:cyclic-jet}
Suppose \(N_0>0\) and \(J_t+{\bf1}_{t\in T}\le N_0\) for all \(t\).
Use \(B\) as in Lemma~\ref{lem:two-differences}, an integral
series \(U\), and the inverse \(g\) of the logarithms \({\cal L}\)
on each matrix entry. Apply the product of all the framed minors
to
\[
          U(g(Z))\exp\!\left(\pm\sum_{i,b}
                        B\big((Z_{sib})_{s\in\mathbb Z/d}\big)\right)
\]
at zero. Each coefficient of the resulting frame polynomial has
valuation at least
\begin{equation}\label{eq:cyclic-gain}
 {\cal F}_p(D)+\sum_{s\in T}(T_s-\gamma W_s)
                    -dN_0N_1(1+\log_p a_*).
\end{equation}
\end{proposition}
\begin{proof}
Use independent indeterminates for the entries of the frames and
the Gauss valuation. At each \(s\in T\) replace the first row
\(Y_{s0}\) by its coordinatewise \(p^m\)-th powers, where
\(m\ge\max_{(s,j)}(-j)\) and \(m\ge0\).
This preserves the valuation of the frame polynomial since
distinct monomials still have distinct exponent vectors. We
construct at \(t\) new rows \(x_{tk}\). The distinguished first
row at \(t\in T\) stays fixed.

To the pair \((s,j)\) ending at
\(t\) assign the covector
\[
                  \lambda_{s,j}=Y_{s0}^{\,p^{m+j}}
\]
using coordinatewise powers of the original first row. Reserve
one distinct nondistinguished row at \(t\) for each of these
covectors. The matrix of their evaluations on the reserved
rows has unit determinant. To see this, reduce the polynomial
modulo the maximal ideal and specialize the reserved rows to
distinct coordinate vectors. These rows are independent of
the covectors, since \(s\ne t\) by parity. In the resulting
determinant each summand is a distinct monomial. At a given
coordinate its variable and exponent recover the pair \((s,j)\)
assigned there. This proves the unit assertion, also when the
list is empty.

Use the inverse evaluation matrix to replace the reserved rows
by dual rows. Subtract the corresponding projections from each
other nondistinguished row. Writing \(r(s,j)\) for the reserved row after this change of basis, the evaluation matrix on all nondistinguished rows becomes
\[
  \lambda_{s,j}(x_{tk})=\begin{cases}1&k=r(s,j),\\0&k\ne r(s,j).\end{cases}
\]
Both basis changes have integral coefficients.
The old minor operators expand in these new rows with integral
coefficients. A determinant with two occurrences of the same
row vanishes.

In any surviving monomial of differentiation with
multi-index \(\kappa\), the distinguished row has degree exactly
\(A_s\), the other rows at \(s\) total \(D_s\), and each of them
has degree at most \(C_s\). At a selected letter, only a minor
of height at least two can contribute to a nondistinguished row;
the row change mixed only those rows with each other. In
particular every coordinate of \(\kappa\) is at most \(a_*\).
The row-by-row bounds on the differentiation orders read
\[
  \sum_b\kappa_{s0b}=A_s,\qquad \sum_{k\ne0,b}\kappa_{skb}=D_s,\qquad \sum_b\kappa_{skb}\le C_s\quad(k\ne0).
\]

In evaluating this monomial make the dual change of variables
\(Z_{sab}=\sum_k x_{sk,a}Z'_{skb}\). It suffices to bound
the series modulo the monomial ideal generated by
\(({Z'}_i^{\kappa_i+1})_i\), since all coefficients entering
\([{Z'}^\kappa]\) survive there. Initially give every variable
weight \(\rho\). By Lemma~\ref{lem:two-differences} only the
exceptional terms need an extra unit of weight for the
exponential. In the expansion of such a term, a mixed choice of
coordinates in either \(p\)-power of a linear form has
multinomial coefficient divisible by \(p\). For the remaining
choices write \(s+\alpha=t+\beta\), with \(s,t\) distinct and of
opposite parity. For coordinates \(i=(s,k,b)\), \(i'=(t,k',b')\)
the bad product is possible only for \(b=b'\). Its coefficient
from that term, divided by the original coefficient of \(B\),
is
\[
                \sum_a x_{sk,a}^{\,p^\alpha}
                           x_{tk',a}^{\,p^\beta}.
\]

If \(i\) is distinguished and \(i'\) is not, let
\(j=\alpha-\beta\). If the product survives truncation then
\((s,j)\) belongs to the list at \(t\). If \(k'\) is not its
reserved row, the last display is zero modulo \(p\). For \(j>0\)
it is congruent to
\((\lambda_{s,j}(x_{tk'}))^{p^\beta}\), and for \(j<0\) it is
congruent to the same evaluation raised to \(p^\beta\), using
\(x_{s0}=Y_{s0}^{p^m}\). Both congruences hold modulo the ideal
\((p)\) of the integral coefficient ring by the binomial
theorem.
In particular, for an unreserved partner the evaluation contributes
\[
  \lambda_{s,j}(x_{tk'})=0\quad\Longrightarrow\quad \sum_a x_{s0,a}^{p^\alpha}x_{tk',a}^{p^\beta}\in(p).
\]
 Thus this case also recovers the missing unit.

We specify extra input weights \(r_i\) for the remaining
products. Give a distinguished row at \(s\) weight increment
\(\epsilon_s\ge0\), requiring \(\epsilon_s\ge p^{-h_s}\) for
finite \(h_s\). Give the nondistinguished rows at \(t\) increment
\(b_t\), raised on the row reserved for \((s,j)\) to
\[
 \begin{cases}
 \max(b_t,1-p^j\epsilon_s),&j>0,\\
 \max(b_t,p^j),&j<0 .
 \end{cases}
\]
For two nondistinguished rows, \eqref{eq:flow-dual} gives
\(p^\alpha r_i+p^\beta r_{i'}\ge1\) after orienting the difference
of exponents so that it is positive.

For two distinguished rows
in that orientation, the first has \(h_s\le\alpha-\beta\), so
the same bound follows from its own increment. For one row of
each kind, the reserved-row formula gives the bound with either
sign of \(j\). Every surviving exponent term now has weight at
least \(\theta\) at input weights \(\rho+r_i\). The exponential
and the integral factor therefore have weight at least zero.
Using the factorial formula as in \eqref{eq:one-letter-jet}, we
obtain the lower bound
\[
       \sum_i\kappa_i(\gamma-r_i)-dN_0N_1(1+\log_p a_*).
\]

For a fixed \(s\in T\) the extra cost of the negative labels
is at most \(W_s\sum_{j>0,\ j\ {\rm odd}}p^{-j}=\gamma W_s\).
For the positive labels and the distinguished row, one can choose
\(\epsilon=\epsilon_s\) so that
\[
 A_s\epsilon+
 \sum_{\substack{j>0\ {\rm odd}\\(s,j)\ {\rm in\ a\ list}}}
       C_{s+j}(1-p^j\epsilon)_+
                \ \le\ A_s\gamma-T_s .
\]
Indeed the right side is the sum of
\(\min(A_sp^{-j},C_{s+j})\) at odd \(j<h_s\), together with
\(A_sp^{-h_s}\) when \(h_s\) is finite. Take the first odd
\(j<h_s\) with \(A_sp^{-j}\le C_{s+j}\), or \(j=h_s\) if none
exists.

At a finite choice set \(\epsilon=p^{-j}\).
For a finite choice the threshold has the two consequences
\[
  A_sp^{-i}>C_{s+i}\quad(0<i<j,\ i\ {\rm odd}),\qquad p^i\epsilon\ge1\quad(i\ge j).
\]
 All earlier
positive-label costs are bounded by their \(C_{s+i}\); later
ones and the one at \(j\) cost zero. The term at \(j\) pays for
\(A_s\epsilon\). If the choice is infinite use \(\epsilon=0\)
and all finite costs are bounded by the respective minima.
This also verifies the lower bound on \(\epsilon_s\). The
nondistinguished baseline is
\(\sum D_t(\gamma-b_t)={\cal F}_p(D)\), and the distinguished
baseline at \(s\), before paying costs, is \(\gamma A_s\).
Adding gives \eqref{eq:cyclic-gain}.
\end{proof}

For comparison with the common chain, write \(V_s=(k-1)w_s\)
with \(k\ge3,\ w_s>0\), and suppose
\[
      D_s=V_s+E_s-2{\bf1}_{s\in T}w_s,\qquad
                         0\le E_s\le4w_s .
\]
Then
\begin{equation}\label{eq:flow-comparison}
            {\cal F}_p(D)\ge {\cal F}_p(V)
                                   -\frac{11}{p}\sum_s w_s .
\end{equation}
To prove this take an optimum \(b\) in \eqref{eq:flow-dual}
for \(V\), and let \(H=\{s:b_s\ge2/p\}\). Complementary
slackness with an optimum flow shows that every vertex of \(H\)
has saturated capacity, and every edge starting there has zero
flow, since its dual constraint is strict. Hence
\(\sum_{s\in H}V_s\) is the total flow into \(H\). Summing the
capacity inequalities shows \(p\sum z_{sj}\le\sum_sV_s\), so

\[
  \sum_H w_s\le p^{-1}\sum_s w_s.
\]
 In the dual formula,
increasing \(V\) by \(E\) decreases its objective at this \(b\)
by at most \(\sum E_s(b_s-\gamma)_+\). On \(H\) use \(b_s\le1\);
outside \(H\) use \(\gamma\ge1/p\). The resulting loss is at
most \((8/p)\sum w_s\). Finally adding \(2{\bf1}_T w\) to
\(D\) increases its optimum dual objective by at most
\(2\gamma\sum_T w_s\), since every feasible \(b\) is
nonnegative. As \(2\gamma\le3/p\), this proves
\eqref{eq:flow-comparison}.
Together the two changes of capacities cost at most
\[
  \frac8p\sum_s w_s+2\gamma\sum_s w_s\le\frac{11}{p}\sum_s w_s.
\]

Taking the smaller numerator \(A_s-A'_s\) and subtracting
\(C_{s+j}\) at every odd \(j\ne h_s\) can only reduce \(T_s\).
Replacing \(D_{s-j}\) by \(H_{s-j}\) in the definition of \(W_s\)
can only increase the loss. If \(H_s\le a\), the number of
positive and negative labels at one destination is at most
\(2\log_p a+2\): for each signed \(j\) there is just one source.
Thus \(2\log_p a+2\le N_0-1\) suffices for the row count; also
\(a_*\le\max(1,a)\). These observations give the constants
\(11\) for the capacity error and \(1\) for the digit loss in
the split bound. A common unit rescaling of the local variables
multiplies a homogeneous jet by a unit and leaves its valuation
unchanged.  \label{cor:cyclic-jets-1}
